% Part: second-order-logic
% Chapter: syntax-and-semantics
% Section: terms-formulas

\documentclass[../../../include/open-logic-section]{subfiles}

\begin{document}

\olfileid{sol}{syn}{frm}

\olsection{ترم‌ها و \printtoken{P}{formula}}

همانند منطق مرتبهٔ اول، عبارت‌های منطق مرتبهٔ دوم از واژگان پایه‌ای
ساخته می‌شوند که \emph{!!{variable}s}، \emph{!!{constant}s}،
\emph{!!{predicate}s} و گاه \emph{!!{function}s} را در بر می‌گیرد.
از این‌ها، همراه با ادات منطقی، سورها و نشانه‌های سجاوندی چون پرانتز و
ویرگول، \emph{ترم‌ها} و \emph{!!{formula}s} ساخته می‌شوند. تفاوت در این
است که منطق مرتبهٔ دوم، افزون بر متغیرهای اشیا، متغیرهایی برای روابط و
توابع نیز دارد و کمیت‌گذاری بر آن‌ها را مجاز می‌شمارد. پس نمادهای منطقی
منطق مرتبهٔ دوم همان نمادهای منطق مرتبهٔ اول‌اند، به‌علاوهٔ موارد زیر:

\begin{enumerate}
\item مجموعه‌ای !!{denumerable} از !!{variable}s رابطه‌ای مرتبهٔ دوم
  با هر آریتی~$n$: $\Obj V_0^n$، $\Obj V_1^n$، $\Obj V_2^n$، \dots
\item مجموعه‌ای !!{denumerable} از !!{variable}s تابعی مرتبهٔ دوم:
  $\Obj u_0^n$، $\Obj u_1^n$، $\Obj u_2^n$، \dots
\end{enumerate}

همان‌گونه که~$x$، $y$ و~$z$ را به‌عنوان فرامتغیرهای متغیرهای مرتبهٔ
اول~$\Obj v_i$ به کار می‌بریم، $X$، $Y$، $Z$ و مانند آن‌ها را
فرامتغیرهای~$\Obj V_i^n$، و~$u$، $v$ و مانند آن‌ها را فرامتغیرهای
$\Obj u_i^n$ به کار خواهیم برد.

\begin{explain}
نمادهای غیرمنطقی یک زبان مرتبهٔ دوم همانند زبان مرتبهٔ اول مشخص
می‌شوند: با فهرست‌کردن !!{constant}s، !!{function}s و !!{predicate}s آن.

در منطق مرتبهٔ اول معمولاً !!{identity} یعنی~$\eq$ گنجانده می‌شود. در
منطق مرتبهٔ اول، نمادهای غیرمنطقیِ زبان~$\Lang{L}$ برای بیان هر مطلب
جالبی ضروری‌اند. البته !!{sentence}sی وجود دارند که هیچ نماد غیرمنطقی
به کار نمی‌برند، اما تنها با~$\eq$ دشوار است مطلب جالبی گفت. در منطق
مرتبهٔ دوم، چون ذخیره‌ای نامحدود از متغیرهای رابطه‌ای و تابعی داریم، حتی
بدون ذخیره‌ای ویژه از نمادهای غیرمنطقی می‌توانیم هر آنچه را در یک زبان
مرتبهٔ اول می‌توان گفت بیان کنیم.
\end{explain}

\begin{defn}[ترم‌های مرتبهٔ دوم]
مجموعهٔ \emph{ترم‌های مرتبهٔ دوم} زبان~$\Lang L$، یعنی~$\TrmSOL[L]$،
با افزودن بند زیر به~\olref[fol][syn][frm]{defn:terms} تعریف می‌شود:
\begin{enumerate}
\item اگر~$u$ یک متغیر تابعی $n$-موضعی و~$t_1$، \dots، $t_n$ ترم
  باشند، آنگاه~$\Atom{u}{t_1, \ldots, t_n}$ یک ترم است.
\end{enumerate}
\end{defn}

\begin{explain}
پس یک ترم مرتبهٔ دوم درست مانند یک ترم مرتبهٔ اول است، جز اینکه در جایی
که ترم مرتبهٔ اول !!a{function} چون~$\Obj{f^n_i}$ دارد، ترم مرتبهٔ دوم
می‌تواند به جای آن متغیر تابعی~$\Obj{u^n_i}$ داشته باشد.
\end{explain}

\begin{defn}[\usetoken{s}{formula} مرتبهٔ دوم]
مجموعهٔ \emph{!!{formula}s مرتبهٔ دوم}~$\FrmSOL[L]$ زبان~$\Lang L$ با
افزودن بندهای زیر به~\olref[fol][syn][frm]{defn:formulas} تعریف می‌شود:
\begin{enumerate}
\item اگر~$X$ یک متغیر محمولی $n$-موضعی و~$t_1$، \dots، $t_n$
  ترم‌های مرتبهٔ دوم زبان~$\Lang L$ باشند، آنگاه
  $\Atom{X}{t_1,\ldots, t_n}$ یک !!{formula} اتمی است.

\tagitem{prvAll}{اگر~$!A$ !!a{formula} و~$u$ متغیری تابعی باشد، آنگاه
  $\lforall[u][!A]$ !!a{formula} است.}{}

\tagitem{prvAll}{اگر~$!A$ !!a{formula} و~$X$ متغیری محمولی باشد، آنگاه
  $\lforall[X][!A]$ !!a{formula} است.}{}

\tagitem{prvEx}{اگر~$!A$ !!a{formula} و~$u$ متغیری تابعی باشد، آنگاه
  $\lexists[u][!A]$ !!a{formula} است.}{}

\tagitem{prvEx}{اگر~$!A$ !!a{formula} و~$X$ متغیری محمولی باشد، آنگاه
  $\lexists[X][!A]$ !!a{formula} است.}{}
\end{enumerate}
\end{defn}

\end{document}
